\section{Mathematical proof of the consistency equivalence}
\label{sec:math}

\subsection{The numerical calculus and exact command interface}

Let $p(a,b)=(a+b)(a+b+1)/2+b$, write $[a;w]=p(a,w)$, and put
\[
 I(a,b)=4p(a,b),\quad M(a,b)=I(a,b)+1,\quad
 U(a,B)=I(B,a)+2,\quad E(a,b)=I(a,b)+3.
\]
Code zero is truth and code one is contradiction. Negation is
$N(B)=I(B,1)$; conjunction and existential quantification are the standard
classical abbreviations. Cantor stacks have total projections:
$\operatorname{hd}(0)=\operatorname{tl}(0)=0$ and $[0;0]=0$.

The unchanged numerical calculus has modus ponens, generalization of
theorems $B\mapsto U(i,B)$, every $\Gamma(B)=I(B,U(B,B))$, and the heads
\begin{align*}
 h_1(D,H)&=(\exists x_D U(D,H))\to H,\\
 h_2(D,A,H)&=U(D,U(A,H))\to U(A,U(D,H)),\\
 h_3(D,A,H)&=M(H,D)\to(E(H,A)\to M(A,D)),\\
 h_4(D,A,H)&=E(H,D)\to(M(A,H)\to M(A,D)),\\
 h_5(D,H)&=\exists x_D E(H,D),\\
 h_6(D,A,H)&=U(D,A)\to(U(D,A\to H)\to U(D,H)),\\
 \rho(D,A,H,d)&=((H\to d)\to D)\to((D\to H)\to(A\to H))\quad(d\le D).
\end{align*}
Its guarded pointed-separation formula is
\begin{align*}
 V(a,b,c)&=M(a,b)\to U(c,M(a,b)),&P(a,b,c)&=c\to M(c,b),\\
 S_*(a,b,c)&=U\bigl(b,(M(a,b)\to P(a,b,c))\wedge(M(c,b)\to c)\bigr),\\
 Q_*(a,b,c)&=U(c,S_*(a,b,c)\to N(P(a,b,c))),\\
 G(a,b,c)&=N(V(a,b,c)\to Q_*(a,b,c)).
\end{align*}
Here $G$ is an axiom when $c\ne a,b$; put $g=G$ under that guard and
$g=0$ otherwise. The code-freshness fact that all labels occurring in a
formula code $B$ are smaller than $B$ justifies $\Gamma(B)$.

A command store is $(t,w,d,h)$, comprising top, stack, offset and cache.
The equations below hold at the inherited command boundaries, where the
scratch register $X$ is zero. They allow arbitrary natural values of the
four displayed store components, not arbitrary unclean physical scratch.
For $b=d+x$, the new source has
\[
 \mathsf{MP}_x(t,w,d,h)=
 \begin{cases}(b,w,b,0),&I(t,b)=h,\\(0,w,b,0),&I(t,b)\ne h,\end{cases}
\]
\[
 S(t,w,d,h)=
 \begin{cases}\mathsf{halt},&h=1,\\
 (h,\operatorname{tl}(w),d,\operatorname{hd}(w)),&h\ne1.\end{cases}
\]
For $D=d+x$, $H=h+y$, axiom generation is
\[
 \mathsf{AX}^{D}_{A,x,y}(t,w,d,h)
 =\bigl(g(H,D,A),[q_1,\ldots,q_9,t;w],0,0\bigr),
\]
where
\[
 (q_1,\ldots,q_9)=
 (\Gamma(D),h_6,h_5,h_4,h_3,h_2,h_1,\rho(D,A,H,d),U(d,h)).
\]
The free-constructor audit checks every displayed formula tree and stack
position. All $q_1,\ldots,q_8,g$ have residue zero modulo four;
$q_9\equiv2\pmod4$. None equals one, without any consistency assumption.

For an explicit distinction, $I(0,2)=20$, and
$\mathsf{MP}_2(0,0,0,20)=(2,0,2,0)$. A subsequent
$\mathsf{AX}^{D}_{0,0,0}$ emits $\Gamma(2)=5712$. The hybrid Gamma-A
parent emits $\Gamma(0)=20$ in this situation, while the cached Gamma-D
parent's MP puts the consequent in $h$ and clears $d$. Thus the new
correlation is not a renaming of the command parameters.

\subsection{A uniform refutation compiler}

Write $A_0=\mathsf{AX}^{D}_{0,0,0}$, and execute words left to right.
Define
\[
 \mathsf{Push}=A_0;S^9;\mathsf{MP}_0,\quad
 \mathsf{Pop}=A_0;S^{11},\quad
 \mathsf{Discard}=\mathsf{MP}_0.
\]
Direct use of the stack equations gives
\begin{align*}
 \mathsf{AX}^{D};S^9(t,w,d,h)&=(q_8,[t;w],0,U(d,h)),\\
 \mathsf{AX}^{D};S^{10}(t,w,d,h)&=(U(d,h),w,0,t),\\
 \mathsf{Push}(t,w,d,h)&=(0,[t;w],0,0),\\
 \mathsf{Discard}(t,w,0,h)&=(0,w,0,0),\\
 \mathsf{Pop}(t,[a;w],d,h)&=
 \begin{cases}\mathsf{halt},&t=1,\\(t,w,0,a),&t\ne1.\end{cases}
\end{align*}
The first two identities are total: their tested caches are zero or
non-one generated heads. Push's final MP fails by the residue difference
between implication and universal codes. Discard follows because both MP
branches agree when the prospective conclusion is zero. These arguments
include the aliases $[0;0]=0$ and $I(0,0)=0$.

For $1\le j\le7$, from $(t,w,0,0)$ the selector
\[
 \mathsf{Discard};\mathsf{AX}^{D}_{A,D,H};S^{j+1};\mathsf{Pop}^{9-j}
\]
returns $(q_j,w,0,0)$. The Gamma-D case is specifically
\[
 \mathsf{Discard};\mathsf{AX}^{D}_{0,B,0};S^2;\mathsf{Pop}^{8},
\]
which selects $\Gamma(B)$ for every $B$. The dominated head, for $u\le a$,
is selected by
\[
 \mathsf{Discard};\mathsf{MP}_u;\mathsf{AX}^{D}_{b,a-u,c};S^9;\mathsf{Pop};
\]
the gate is selected by
\[
 \mathsf{Discard};\mathsf{AX}^{D}_{A,D,H};\mathsf{Pop}^{10}.
\]
Generalization is
\[
 \mathsf{Gen}_i=\mathsf{Push};\mathsf{MP}_i;S;
                  \mathsf{AX}^{D}_{i,0,0};S^{10}.
\]
It sends $(a,w,0,0)$ to $(U(i,a),w,0,0)$ even when $a=1$.

For each numerical derivation $p$, construct a word $C(p)$ satisfying:
from every $(t,w,0,0)$ it either halts or returns
$(\operatorname{concl}(p),w,0,0)$. Use the displayed selectors for axioms,
and append $\mathsf{Gen}_i$ at generalization. For MP with premise
proofs $p_0:a$ and $p_1:I(a,b)$ use
\[
 \mathsf{Discard};\mathsf{Push};C(p_1);\mathsf{Push};C(p_0);
 \mathsf{Push};S;S;\mathsf{MP}_b;\mathsf{Pop}.
\]
The reserved zero gives $(0,[a,I(a,b),0;w],0,0)$ before the two Shifts.
The second Shift halts if $a=1$; otherwise MP returns
$(b,[0;w],b,0)$. The final Pop halts if $b=1$; otherwise it consumes the
reserved zero and returns $(b,w,0,0)$. This proves the invariant by
induction on the finite proof tree. In particular a proof of one yields
$A_0;C(p);\mathsf{Pop}$, which halts from every incoming store.
The all-store qualification is essential because the scheduler retains
its store between attempted words.

The unchanged scheduler enumerates all finite command words with positive
continuation codes. Specifically, for suffix code $z$,
\[
 e(S,z)=p(0,z),\quad e(\mathsf{MP}_x,z)=p(1,p(x,z)),\quad
 e(\mathsf{AX}_{A,x,y},z)=p(A+2,p(y,p(x,z))).
\]
The list encoder uses $E(\epsilon)=0$, $E(c::w)=1+e(c,E(w))$.
Positive continuations decrease, so each nonhalting candidate finishes;
fresh candidates run through consecutive integers from two. Prefixing
the universally halting word by arbitrary AX commands gives unbounded
fresh candidate codes and avoids the omitted initial values.
All proof-word constructors, bounded candidate traversals and finite
trace concatenations are primitive recursive, and PA verifies their
correctness by ordinary induction.

Conversely, annotate $t$, $h$ and every canonical stack entry with a
numerical-calculus proof. Zero carries its fixed truth proof; numerical
zero-stack padding is treated explicitly. AX emits axioms, a generalization
of the old cache, and the old top. A successful MP produces its conclusion;
a failed MP produces truth. Shift transports proofs and halts only on a
cached proof of one. The offset needs no proof annotation. This is a
syntactic invariant, not an appeal to arithmetic reflection.

\subsection{The two inserted zero tests and exact physical realization}

The command argument concerns the ordinary source with unconditional
returns. The literal implementation replaces its site-26 jump by a
zero test on $X$ and its site-3553 final AX jump by a zero test on $D$.
The positivity abstraction records for each of the nine registers whether
it is zero or positive. Its recomputed least closure has 263,236
location/mask pairs. Both ordinary and specialized programs have this
same closed set and identical successor sets there; the tested register
is zero in every mask at each changed site. The finite checker verifies
these facts. Induction on executions therefore proves exact equality
of the two seeded source runs. Reusing the older Gamma-D closure would
have been unjustified: its 263,244-pair set is different.

The new source has 3,650 sites and 4,752 decrement-branch-inclusive
outcomes. The independent literal checker reconstructs its sites from
the pinned Gamma-D source by explicit address splices, compares two
separate control-flow resolvers, rederives four tape-list lemmas from the
literal table, and proves every source outcome for arbitrary natural
register values, arbitrary natural header gap and every nonempty finite
spare list. It uses no arithmetic-box oracle. The ten simultaneous
history folds pass forty local two-cell contracts with arbitrary exterior
tape. Every source simulation lasts strictly positive time and preserves
the complete frame, including spare-list effects.

Exact blank execution reaches native state 84, head $-2$, source site 25,
logical registers $(2,0,0,0,0,0,0,0,0)$ and 26 spare registers of value one
after 58,014 transitions, with no earlier halt. Independent set and
bytearray simulations compare the entire tape and zero exterior.
Concatenating positive simulations transports source halts to tape halts.
If the tape first halted after $T$ transitions while the first $T+1$
source steps were nonhalting, their simulations would produce the same
deterministic physical run for more than $T$ nonhalting transitions,
a contradiction. This supplies both finite-trace implications in PA.

\subsection{Consistency conclusion and proof boundary}

The supplied separate-head calibration provides primitive-recursive
translations between PA refutations and refutations of the numerical
calculus, with correctness provable in PA. It uses diagonal-code freshness,
reverse membership rewriting, neighborhood equality, and arithmetic
finite-comprehension proof templates. The present construction retains
that calculus exactly; Gamma-D changes only the correlated command
parameters used to select its already available diagonal instances.
Composing that calibration, the new command compiler and annotation
invariant, the fresh zero-test proof and the literal simulations gives
\[
 \mathrm{PA}\vdash
 \bigl(\neg\operatorname{Halt}(\mathrm{PA194DHS},\mathrm{blank})
             \leftrightarrow\operatorname{Con}(\mathrm{PA})\bigr).
\]
This is an ordinary mathematical PA-provability argument with executable
finite and universal symbolic certificates. The inherited logical
calibration is an explicit dependency; these certificates are not a
new end-to-end Lean or Rocq theorem. Standard-model nonhalting remains
conditional on the consistency of PA.
